% ══════════════════════════════════════════════════════════════════════
% SECTION VIII
% ══════════════════════════════════════════════════════════════════════
\section{Future Directions: Closing the Measurement Imbalance}
\label{sec:future}

\subsection{Measuring the Gap: A Cross-Benchmark Coverage Analysis}

The measurement imbalance identified throughout this survey is not only visible in aggregate statistics---it is legible benchmark by benchmark. Table~\ref{tab:benchmark_coverage} maps nine representative benchmarks from the 2024--2026 literature against the four evaluation pillars and three operational dimensions (economic, adherence, resilience) drawn from the CLEAR framework. Ratings reflect the authors' synthesis of published evaluation protocols and should be read as characterizing what each benchmark \emph{explicitly evaluates}, not as an overall quality assessment.

\begin{table*}[htbp]
\centering
\caption{Cross-benchmark coverage analysis. Ratings: $\checkmark$\,=\,explicitly evaluated, $\circ$\,=\,partially or implicitly covered, ---\,=\,not addressed. Assessments are the authors' synthesis based on published evaluation protocols; official benchmark descriptions may differ.}
\label{tab:benchmark_coverage}
\footnotesize
\setlength{\tabcolsep}{4pt}
\resizebox{\textwidth}{!}{%
\begin{tabular}{@{}lccccccc@{}}
\toprule
\textbf{Benchmark} & \textbf{Core Model} & \textbf{Memory} & \textbf{Tool Use} & \textbf{Environment} & \textbf{Economic (Cost/Lat.)} & \textbf{Adherence} & \textbf{Resilience} \\
\midrule
SWE-bench~\citep{beyond_benchmark_islands}    & $\checkmark$ & $\circ$ & $\checkmark$ & $\circ$    & ---        & ---        & ---        \\
WebArena~\citep{beyond_benchmark_islands}     & $\checkmark$ & ---     & $\checkmark$ & $\checkmark$ & ---        & ---        & ---        \\
AgentBench~\citep{comprehensive_empirical}    & $\checkmark$ & ---     & $\checkmark$ & $\checkmark$ & ---        & ---        & ---        \\
GAIA~\citep{beyond_benchmark_islands}         & $\checkmark$ & ---     & $\checkmark$ & $\circ$    & ---        & ---        & ---        \\
$\tau$-bench~\citep{beyond_benchmark_islands} & $\checkmark$ & ---     & $\checkmark$ & $\circ$    & ---        & $\circ$    & ---        \\
HAAF~\citep{beyond_benchmark_islands}         & $\checkmark$ & $\circ$ & $\checkmark$ & $\checkmark$ & $\circ$    & $\checkmark$ & $\circ$    \\
AgentDrive~\citep{beyond_benchmark_islands}   & $\checkmark$ & $\circ$ & $\circ$    & $\checkmark$ & ---        & ---        & $\checkmark$ \\
Agent GPA~\citep{agent_gpa}                   & $\checkmark$ & $\circ$ & $\circ$    & ---        & ---        & ---        & ---        \\
Mind2Web~\citep{beyond_benchmark_islands}     & $\checkmark$ & ---     & $\checkmark$ & $\circ$    & ---        & ---        & ---        \\
\midrule
\textbf{Coverage (\%)} & \textbf{100\%} & \textbf{approx. 33\%} & \textbf{approx. 78\%} & \textbf{approx. 67\%} & \textbf{approx. 11\%} & \textbf{approx. 22\%} & \textbf{approx. 22\%} \\
\bottomrule
\end{tabular}
}
\end{table*}

The pattern is stark. Core Model coverage is universal; Tool Use and Environment are reasonably well-represented. Memory evaluation is partial in only one-third of benchmarks and absent in the rest. Economic, Adherence, and Resilience coverage each falls below 25\% across the set. HAAF is the only benchmark in this survey that explicitly addresses four or more of the seven dimensions. This table is not an argument that every benchmark must cover every dimension; domain-specific benchmarks rightly focus on the dimensions most relevant to their deployment context. It is, however, a concrete illustration of where the research community has concentrated its measurement effort and where the gaps are largest---gaps that Table~\ref{tab:benchmark_coverage} makes difficult to dispute.

\subsection{Integrating Human-Centered and Economic Metrics}

A systematic review of papers from 2023 to 2025 conducted by \citet{measurement_imbalance} reveals a \keyword{measurement imbalance} within that study's review corpus: approximately 83\% of the evaluations analyzed focused primarily on technical performance, while only around 30\% addressed human-centered factors and a similar proportion considered economic impacts. These percentages are specific to the coding and review criteria used by that study and should be interpreted as characterizing a particular sample of the literature, not as universal constants. The directional finding---that technical metrics dominate over human-centered and economic ones---is consistent with the broader pattern visible across the surveyed literature. Figure~\ref{fig:measurement_imbalance} visualizes this imbalance.

\begin{figure}[htbp]
\centering
% Measurement Imbalance in Agentic AI Evaluation
\resizebox{0.68\columnwidth}{!}{%
\begin{tikzpicture}
\begin{axis}[
    ybar,
    width=4.9cm,
    height=3.5cm,
    bar width=7pt,
    ymin=0,
    ymax=90,
    ylabel={Share of reviewed studies (\%)},
    ylabel style={font=\tiny, scale=0.72, transform shape, at={(axis description cs:-0.18,0.5)}},
    symbolic x coords={Technical, Human, Economic, Security},
    xtick=data,
    x tick label style={font=\tiny, scale=0.72, transform shape, rotate=18, anchor=east},
    ytick={0,20,40,60,80},
    yticklabel style={font=\tiny, scale=0.72, transform shape},
    nodes near coords,
    nodes near coords style={font=\tiny\bfseries, scale=0.72, transform shape},
    every axis plot/.append style={draw=none, fill opacity=0.9},
    grid=major,
    grid style={gray!15},
    axis lines*=left,
    enlarge x limits=0.08,
    clip=false
]

\addplot[fill=alcheblue] coordinates {
    (Technical, 83)
    (Human, 30)
    (Economic, 30)
    (Security, 12)
};

%\node[font=\scriptsize, text=riskred, align=center] at (axis cs:Security,20)
%    {Technical emphasis remains dominant;\\security coverage is thinnest.};

\end{axis}
\end{tikzpicture}
}
\vspace{-0.2em}
\fontsize{6}{6.8}\selectfont
\caption{The measurement imbalance in agentic AI evaluation. Technical metrics dominate the literature, while security and human-centered considerations remain critically underrepresented.}
\label{fig:measurement_imbalance}
\end{figure}

The measurement imbalance is not merely an academic concern. It has direct business consequences:

\begin{itemize}
    \item Organizations invest in agents that perform well on technical benchmarks but fail to deliver value to users.
    \item Procurement decisions are made based on accuracy metrics that do not account for cost, latency, or reliability.
    \item Agents are deployed that are technically capable but operationally fragile, leading to production incidents and erosion of trust.
\end{itemize}

Closing this imbalance requires three shifts:

\begin{enumerate}
    \item \textbf{Mandatory economic reporting}: Every benchmark result should include CNA or an equivalent cost-normalized metric. Accuracy without cost context is misleading.
    \item \textbf{Human-centered evaluation integration}: User satisfaction, task completion from the user's perspective, and the quality of the human-agent interaction should be standard evaluation dimensions, not afterthoughts.
    \item \textbf{Multi-dimensional deployment criteria}: The CLEAR framework (or equivalent) should be adopted as a standard for deployment readiness, requiring assessment across all five dimensions.
\end{enumerate}

\subsection{The Role of Standardized Web Conduct and Multi-Agent Interaction Protocols}

As agentic AI matures, the industry is moving toward a more nuanced understanding of ``agenticness'' as a spectrum rather than a binary property~\citep{aiagentindex}. Systems are increasingly being designed with a \keyword{flexible thinking budget}---skipping complex orchestrators for reactive tasks and allocating more tokens and time for deep reasoning tasks like supply chain planning.

This evolution creates new evaluation challenges:

\begin{itemize}
    \item \textbf{Multi-agent evaluation}: When multiple agents interact, emergent behaviors can arise that are not predictable from evaluating each agent in isolation. Metrics for coordination, communication efficiency, and collective decision-making are needed.
    \item \textbf{Standardized interaction protocols}: As agents interact with web services, APIs, and each other, standardized protocols for agent conduct are needed---analogous to robots.txt for web crawlers, but for autonomous AI agents.
    \item \textbf{Cross-platform benchmarking}: Current benchmarks are fragmented and non-comparable. Industry-standard evaluation protocols are needed to enable meaningful comparison across agents and platforms.
\end{itemize}

The development of these standards will require collaboration across industry, academia, and regulatory bodies. Organizations that contribute to standardization efforts will have a voice in shaping the evaluation landscape.

\subsection{Toward Verifiable and Trustworthy Autonomous Agents}

Industry practitioners have articulated aspirational targets for enterprise-grade agentic systems---including 95\%+ task accuracy in multi-turn reasoning workflows~\citep{mckinsey_agentic_ai}---though such figures reflect industry planning horizons and marketing commitments rather than academic consensus benchmarks. Whether 95\% represents a meaningful or achievable threshold depends heavily on task definition, evaluation methodology, and deployment context. We cite this figure as an indicator of where the practitioner community is directing investment, not as a validated scientific target. Reaching even a well-defined accuracy target requires closing the measurement imbalance---integrating technical performance, human-centered assessment, and economic efficiency into a coherent evaluation approach.

The methodological building blocks now exist:

\begin{itemize}
    \item \textbf{Trajectory metrics} that evaluate the complete execution path, not just the final output.
    \item \textbf{Agent-as-a-Judge} paradigms that provide scalable, human-aligned evaluation.
    \item \textbf{CNA} and the CLEAR framework that integrate economic and operational dimensions.
    \item \textbf{Simulation-based pre-deployment testing} that surfaces issues before production.
    \item \textbf{The Four-Pillar Model} that decomposes agent assessment into testable components.
    \item \textbf{The HAAF scenario manifold} that weights evaluation by deployment risk.
\end{itemize}

The challenge is adoption. Most teams are still optimizing for the metric they can measure most easily, not the metric that matters most for production.


\subsection{The Evaluation Horizon Problem}

A less-discussed structural challenge for the field is what we term the \keyword{evaluation horizon}---the finite time interval over which a benchmark retains its discriminative validity before degradation through training data contamination.

As agentic models are continuously trained on large, internet-scale corpora, published benchmarks become part of the training distribution over time. A benchmark released in 2024 may appear on forum discussions, code repositories, and synthetic data generators used in post-training runs in 2025 and 2026. The result is that benchmark performance increasingly reflects how much of the benchmark's specific task distribution an agent has been trained on---a confound that worsens with each model generation and is difficult to detect from evaluation scores alone.

We define the evaluation horizon $H(B, \epsilon)$ of benchmark $B$ as the time interval after which the benchmark's discriminative validity---its ability to distinguish better agents from worse ones---drops below threshold $\epsilon$. If $\delta(t)$ denotes the contamination-adjusted discriminative validity of $B$ at time $t$ after release $t_0$:
\begin{equation}
    H(B, \epsilon) = \inf \{ t > t_0 : \delta(t) < \epsilon \}
\end{equation}
For widely-cited, publicly released benchmarks, this horizon may be as short as 6--18 months given the pace of model releases. This framing has direct implications: benchmark maintainers should budget for regular task refresh cycles; research teams should prefer held-out evaluation sets not released publicly; and the community should weight recently-released, unpublished benchmarks more heavily in comparative assessments. The evaluation horizon problem is also a structural argument for simulation-based and generative evaluation methodologies---task distributions that are generated at evaluation time are harder to contaminate than fixed task sets.

% ══════════════════════════════════════════════════════════════════════
% SECTION IX
% ══════════════════════════════════════════════════════════════════════
\section{Conclusion and Strategic Recommendations}
\label{sec:conclusion}

This survey has synthesized the agentic AI evaluation literature of 2024--2026, identifying structural gaps, reviewing proposed frameworks, and critically examining the tooling ecosystem. The key finding is that evaluation methodologies inherited from static NLP are systematically inadequate for the non-deterministic, multi-step, tool-mediated behavior of agentic systems. Task completion is a necessary but insufficient condition for production readiness. The CLEAR dimensions and Cost-Normalized Accuracy proposed here are conceptual synthesis contributions intended to prompt community standardization discussion---not empirically validated instruments ready for direct adoption. Similarly, the Four-Pillar Model and success masking framing are organizational syntheses rather than novel empirical discoveries.

Organizations and researchers must prioritize trajectory-level observability to understand why agents succeed or fail, particularly in multi-turn interactions where compounding errors are common.

The adoption of Agent-as-a-Judge paradigms offers a scalable path to achieving human-level evaluation quality while significantly reducing costs. However, this must be balanced with rigorous stress-testing of those judges using libraries like the Judge Reliability Harness to ensure they are not merely reflecting their own internal biases.

The economic dimension of agency---specifically the trade-off between accuracy and token consumption---must become a central focus of both academic research and enterprise procurement to prevent the explosion of inference bills and the deployment of operationally fragile systems.

\subsection{Strategic Recommendations}

\subsubsection{For Research Teams}

\begin{enumerate}
    \item \textbf{Report CNA alongside accuracy.} Any benchmark result reported without cost normalization is incomplete and potentially misleading.
    \item \textbf{Adopt trajectory-level evaluation.} Move beyond outcome-only metrics to assess the complete execution path.
    \item \textbf{Stress-test your judges.} Implement JRH-style perturbation testing for any LLM-based evaluation pipeline.
    \item \textbf{Close the measurement imbalance.} Integrate human-centered and economic metrics into evaluation frameworks alongside technical performance.
\end{enumerate}

\subsubsection{For Engineering Teams}

\begin{enumerate}
    \item \textbf{Instrument your traces before you instrument your benchmarks.} Trajectory observability is a prerequisite for meaningful evaluation. If you cannot replay the reasoning path for a failed task, you cannot diagnose the failure.
    \item \textbf{Use multi-turn simulation to find issues before users do.} Tools like ArkSim can surface context loss, idempotency failures, and unexpected conversation paths before they appear in production logs.
    \item \textbf{Apply the four pillars selectively by risk.} Not every workflow requires evaluation across all four pillars at equal depth. Triage by consequence: high-stakes workflows warrant full trajectory evaluation; lower-stakes workflows can tolerate lighter-weight assessment.
    \item \textbf{Implement the CLEAR framework for deployment governance.} No agent should be promoted to production without assessment across all five CLEAR dimensions.
\end{enumerate}

\subsubsection{For Leadership and Procurement}

\begin{enumerate}
    \item \textbf{Require multi-dimensional evaluation in vendor assessments.} Accuracy-only benchmarks are insufficient for procurement decisions. Demand CNA, CLEAR assessments, and trajectory-level evaluation results.
    \item \textbf{Budget for evaluation infrastructure.} The cost of inadequate evaluation---production incidents, compliance violations, eroded trust---far exceeds the cost of building proper evaluation infrastructure.
    \item \textbf{Invest in organizational learning.} Evaluation data should inform agent design, deployment decisions, and risk calibration across the organization.
\end{enumerate}

\vspace{1cm}
\noindent By adopting a multi-dimensional approach that spans technical, human-centered, and economic factors, the industry can bridge the gap between impressive laboratory results and reliable, production-grade autonomous infrastructure.

%\vspace{2cm}
%\begin{center}
%\rule{0.3\textwidth}{0.4pt}\\[6pt]
%\textit{\small \textcopyright\ 2026 Alchedata. All rights reserved.}
%\end{center}